% Part: incompleteness
% Chapter: representability-in-q
% Section: basic-representable

\documentclass[../../../include/open-logic-section]{subfiles}

\begin{document}

\olfileid{inc}{req}{bre}
\olsection{मूलभूत फलने~$\Th{Q}$ मध्ये निरूपणीय आहेत}

सर्वप्रथम सर्व मूलभूत फलने~$\Th{Q}$ मध्ये निरूपणीय आहेत, हे
दाखवायचे आहे. अखेरीस, $k$ आदाने असलेल्या प्रत्येक मूलभूत
फलनासाठी $f(x_0,\dots,x_{k-1})$ त्याचे निरूपण करणारे
!!a{formula} सूत्र $!A_f(x_0,\dots,x_{k-1},y)$ कसे द्यायचे,
हे दाखवावे लागेल.

शून्य, उत्तरवर्ती, बेरीज, गुणाकार, समतेचे जातिबोधक फल आणि
प्रक्षेपण फलने यांची निरूपणीयता दाखवता येईल. प्रत्येक बाबतीत
योग्य निरूपक सूत्र सहज सुचते; उदाहरणार्थ, शून्याचे निरूपण
$y = \Obj 0$ या सूत्राने, उत्तरवर्ती फलनाचे !!{formula}
$x_0' = y$ या सूत्राने आणि बेरीजेचे !!{formula}
$(x_0 + x_1) = y$ या सूत्राने होते. खरे काम म्हणजे संबंधित
!!{sentence}s $\Th{Q}$ मध्ये सिद्ध होतात, हे दाखवणे. उदाहरणार्थ,
वरील !!{formula} सूत्र बेरीजेचे निरूपण करते असे म्हणण्यासाठी
नैसर्गिक संख्यांच्या प्रत्येक जोडी $m$ आणि $n$ साठी $\Th{Q}$
पुढील विधाने सिद्ध करते, हे दाखवावे लागते:
\begin{align*}
& \eq[\num n + \num m][\num {n+m}] \text{ आणि}\\
& \lforall[y][(\eq[(\num n + \num m)][y] \lif \eq[y][\num{n+m}])].
\end{align*}

\begin{prop}
\ollabel{prop:rep-zero}
शून्य फलन $\Zero(x) = 0$ चे~$\Th{Q}$ मध्ये
$!A_{\Zero}(x,y) \ident \eq[y][\Obj 0]$ या सूत्राने निरूपण होते.
\end{prop}

\begin{prop}
\ollabel{prop:rep-succ}
उत्तरवर्ती फलन $\Succ(x) = x+1$ चे~$\Th{Q}$ मध्ये
$!A_{\Succ}(x,y) \ident \eq[y][x']$ या सूत्राने निरूपण होते.
\end{prop}

\begin{prop}
\ollabel{prop:rep-proj}
प्रक्षेपण फलन $\Proj{n}{i}(x_0, \dots, x_{n-1}) = x_i$ चे
~$\Th{Q}$ मध्ये पुढील सूत्राने निरूपण होते:
\[!A_{\Proj{n}{i}}(x_0, \dots, x_{n-1}, y) \ident \eq[y][x_i].\]
\end{prop}

\begin{prob}
$\eq[y][\Obj 0]$, $\eq[y][x']$ आणि $\eq[y][x_i]$
ही सूत्रे अनुक्रमे $\Zero$, $\Succ$ आणि $\Proj{n}{i}$ यांचे
निरूपण करतात, हे सिद्ध करा.
\end{prob}

\begin{prop}
\ollabel{prop:rep-id}
समता~$=$ हिचे जातिबोधक फल,
\[
\Char{=}(x_0, x_1) =
\begin{cases}
  1 & \text{जर } x_0 =x_1\\
  0 & \text{अन्यथा}
\end{cases}
\]
याचे~$\Th{Q}$ मध्ये पुढील सूत्राने निरूपण होते:
\[
  !A_{\Char{=}}(x_0, x_1, y) \ident (\eq[x_0][x_1] \land \eq[y][\num{1}]) \lor (\eq/[x_0][x_1] \land
\eq[y][\num{0}]).
\]
\end{prop}

या पुराव्यासाठी पुढील साहाय्यक विधान लागते.

\begin{lem}
\ollabel{lem:q-proves-neq} $n$ आणि $m$ या नैसर्गिक संख्या असू द्या.
$n \neq m$ असेल, तर $\Th{Q} \Proves \eq/[\num n][\num m]$.
\end{lem}

\begin{proof}
$n$ वर विगमन वापरून, प्रत्येक $m$ साठी $n \neq m$ असेल तर
$Q \Proves \eq/[\num n][\num m]$, हे दाखवू.

आधार टप्प्यात $n = 0$. जर $m$ हे $0$ च्या बरोबर नसेल,
तर एखादी नैसर्गिक संख्या $k$ अशी असते की $m = k +
1$. $\lforall[x][\eq/[0][x']]$ असे सांगणारे स्वयंसिद्धक
आपल्याकडे आहे. संख्यकाच्या एका स्वयंसिद्धकानुसार $x$ च्या जागी
$\num k$ ठेवून $\eq/[0][\num k']$ हा निष्कर्ष मिळतो. पण
$\num k'$ म्हणजेच $\num m$.

विगमन टप्प्यात, $n$ साठी विधान खरे आहे असे धरून $n+1$ चा
विचार करू. $m$ ही कोणतीही नैसर्गिक संख्या असू द्या. दोन शक्यता
आहेत: $m = 0$, किंवा एखाद्या $k$ साठी $m = k+1$. पहिली
शक्यता वरप्रमाणे हाताळता येते. दुसऱ्या शक्यतेत $n+1 \neq
k+1$ असे समजा. मग $n \neq k$. $n$ साठीच्या विगमन गृहीतकाने
$\Th{Q} \Proves \eq/[\num n][\num k]$. आपल्याकडे
$\lforall[x][\lforall[y][\eq[x'][y'] \lif \eq[x][y]]]$
असे सांगणारे स्वयंसिद्धक आहे. संख्यकाच्या एका स्वयंसिद्धकाने
$\eq[\num n'][\num k'] \lif \eq[\num n][\num
  k]$ मिळते. विधानीय तर्कशास्त्र वापरून $\Th{Q}$ मध्ये
$\eq/[\num n][\num k] \lif \eq/[\num n'][\num k']$
हा निष्कर्ष मिळतो. आता अभिव्यंजन-निर्मूलन नियमाने
$\eq/[\num n'][\num k']$ मिळते; आणि $\num k'$ म्हणजेच
$\num m$ असल्याने आपल्याला हेच दाखवायचे होते.
\end{proof}

\begin{explain}
हे साहाय्यक विधान फार मोठी गोष्ट सांगत नाही: वेगवेगळे संख्यांक
वेगवेगळ्या वस्तू दर्शवतात, हे $\Th{Q}$ सिद्ध करू शकते,
असे त्याचे सार आहे. उदाहरणार्थ, $\Th{Q}$ हे
$0'' \neq 0'''$ सिद्ध करते. पण हे सर्वसाधारणपणे खरे आहे
हे दाखवताना काही काळजी घ्यावी लागते. येथे विगमन वापरले असले,
तरी ते $\Th{Q}$ च्या \emph{बाहेर} केलेले विगमन आहे, हेही लक्षात घ्या.
\end{explain}

\begin{proof}[\olref{prop:rep-id} चा पुरावा]
$n = m$ असेल, तर $\num{n}$ आणि $\num{m}$ ही एकच पदे आहेत
आणि $\Char{=}(n, m) = 1$. तसेच $\Th{Q} \Proves (\eq[\num{n}][\num{m}] \land
\eq[\num{1}][\num{1}])$, म्हणून ते $!A_=(\num{n}, \num{m},
\num{1})$ सिद्ध करते. $n \neq m$ असेल, तर $\Char=(n, m) = 0$.
\olref{lem:q-proves-neq} नुसार $\Th{Q} \Proves \eq/[\num{n}][\num{m}]$;
म्हणून $(\eq/[\num{n}][\num{m}] \land \Obj 0 = \Obj 0)$
हेही सिद्ध होते. त्यामुळे $\Th{Q}
\Proves !A_=(\num{n}, \num{m}, \num{0})$.

दुसऱ्या अटीसाठीही दोन शक्यता आहेत. $n = m$ असेल, तर
$\Th{Q} \Proves \lforall[y][(!A_=(\num{n}, \num{m}, y)
  \lif \eq[y][\num{1}])]$ हे दाखवायचे आहे. अनौपचारिकपणे
युक्तिवाद करू: $!A_=(\num{n}, \num{m}, y)$, म्हणजेच
\[
(\eq[\num{n}][\num{n}] \land \eq[y][\num{1}]) \lor
(\eq/[\num{n}][\num{n}] \land \eq[y][\num{0}])
\]
असे समजा. डाव्या विकल्पातून तर्कशास्त्राने
$\eq[y][\num{1}]$ मिळते. उजवा विकल्प मात्र
तर्कशास्त्राने सिद्ध होणाऱ्या $\eq[\num{n}][\num{n}]$
शी विसंगत आहे.

दुसरीकडे, $n \neq m$ असेल, तर $!A_=(\num{n},
\num{m}, y)$ असे आहे:
\[
(\eq[\num{n}][\num{m}] \land \eq[y][\num{1}]) \lor
(\eq/[\num{n}][\num{m}] \land \eq[y][\num{0}])
\]
येथे डावा विकल्प \olref{lem:q-proves-neq} नुसार
$\Th{Q}$ मध्ये सिद्ध होणाऱ्या $\eq/[\num{n}][\num{m}]$
शी विसंगत आहे; उजव्या विकल्पातून $\eq[y][\num{0}]$ मिळते.
\end{proof}

\begin{prop}
\ollabel{prop:rep-add}
बेरीज फलन $\Add(x_0, x_1) = x_0+x_1$ चे~$\Th{Q}$ मध्ये
पुढील सूत्राने निरूपण होते:
\[
  !A_{\Add}(x_0, x_1, y) \ident \eq[y][(x_0 + x_1)].
\]
\end{prop}

\begin{lem}
\ollabel{lem:q-proves-add}
$\Th{Q} \Proves \eq[(\num{n} + \num{m})][\num{n+m}]$
\end{lem}

\begin{proof}
$m$ वर विगमन वापरून हे सिद्ध करू. $m = 0$ असेल, तर
$\Th{Q} \Proves \eq[(\num{n} + \Obj 0)][\num{n}]$
हे विधान आहे. ते स्वयंसिद्धक~$!Q_4$ वरून मिळते. आता
$m$ साठी विधान खरे आहे असे समजून $m+1$ साठी, म्हणजे
$\Th{Q} \Proves \eq[(\num{n} +
  \num{m+1})][\num{n+m+1}]$ हे सिद्ध करू. लक्षात घ्या,
$\num{m+1}$ म्हणजेच $\num{m}'$ आणि $\num{n+m+1}$ म्हणजेच
$\num{n+m}'$. स्वयंसिद्धक $!Q_5$ ने $\Th{Q}
\Proves \eq[(\num{n} + \num{m}')][(\num{n}+\num{m})']$.
विगमन गृहीतकाने $\Th{Q} \Proves \eq[(\num{n} + \num{m})][\num{n+m}]$.
म्हणून $\Th{Q} \Proves \eq[(\num{n} + \num{m}')][\num{n+m}']$.
\end{proof}

\begin{proof}[\olref{prop:rep-add} चा पुरावा]
$\Add$ चे निरूपण करणारे !!{formula} सूत्र~$!A_\Add(x_0, x_1, y)$
म्हणजे $\eq[y][(x_0 + x_1)]$. प्रथम, $\Add(n, m) = k$
असेल, तर $\Th{Q} \Proves !A_\Add(\num{n}, \num{m}, \num{k})$,
म्हणजे $\Th{Q}
\Proves \eq[\num{k}][(\num{n} + \num{m})]$, हे दाखवू.
पण $k = n + m$ असल्यामुळे $\num{k}$ म्हणजेच $\num{n+m}$;
आणि \olref{lem:q-proves-add} मध्ये आपण
$\Th{Q} \Proves \eq[(\num{n} +
  \num{m})][\num{n+m}]$ हे सिद्ध केले आहे. समतेच्या
सममितीने त्यावरून अपेक्षित उलट क्रमातील समता मिळते.

$\Add(n, m) = k$ असेल, तर पुढील विधानही दाखवायचे आहे:
\[
\Th{Q} \Proves \lforall[y][(!A_\Add(\num{n}, \num{m}, y) \lif
  \eq[y][\num{k}])].
\]
निरूपक सूत्र गृहीत धरून, समतेच्या सममितीने
$\eq[(\num{n} + \num{m})][y]$ असे समजा. कारण
\[
\Th{Q} \Proves \eq[(\num{n}+\num{m})][\num{n+m}],
\]
म्हणून डावीकडील पदाऐवजी $\num{n+m}$ ठेवून कोणत्याही~$y$
साठी $\eq[\num{n+m}][y]$ मिळते. समतेच्या सममितीने
अपेक्षित क्रमातील समता मिळते; उलट अभिव्यंजनही याच सिद्ध
समतेवरून आणि समतेच्या नियमांवरून मिळते.
\end{proof}

\begin{prop}
\ollabel{prop:rep-mult}
गुणाकार फलन $\Mult(x_0, x_1) = x_0 \cdot x_1$ चे
~$\Th{Q}$ मध्ये पुढील सूत्राने निरूपण होते:
\[
  !A_{\Mult}(x_0, x_1, y) \ident y = (x_0 \times x_1).
\]
\end{prop}

\begin{proof} सरावासाठी. \end{proof}

\begin{lem}
\ollabel{lem:q-proves-mult}
$\Th{Q} \Proves \eq[(\num{n} \times \num{m})][\num{n \cdot m}]$
\end{lem}

\begin{proof} सरावासाठी. \end{proof}

\begin{prob}
\olref[inc][req][bre]{lem:q-proves-mult} सिद्ध करा.
\end{prob}

\begin{prob}
\olref[inc][req][bre]{lem:q-proves-mult} वापरून
\olref[inc][req][bre]{prop:rep-mult} सिद्ध करा.
\end{prob}

\begin{explain}
  आठवून पाहा: अंकगणिताच्या भाषेतील फलनचिन्हासाठी आपण $\times$
  वापरतो, तर संख्यांवरील नेहमीच्या गुणाकारासाठी $\cdot$ वापरतो.
  म्हणून संख्यादर्शक व्यंजकांमध्ये (उदाहरणार्थ, $m \cdot n$)
  $\cdot$ येऊ शकते; पण अंकगणिताच्या भाषेतील पदांमध्येच
  (उदाहरणार्थ, $(\num{m} \times
  \num{n})$) $\times$ येते. आणखी गोंधळाची बाब म्हणजे $+$
  हे चिन्ह !!{function} आणि बेरीजक्रिया या दोन्हींसाठी वापरतो.
  ते पदांच्या मध्ये आल्यास---उदाहरणार्थ, $(\num{n} + \num{m})$
  मध्ये---अंकगणिताच्या भाषेतील $2$ आदानांचे !!{function} असते;
  आणि संख्यांच्या मध्ये आल्यास---उदाहरणार्थ, $n+m$ मध्ये---
  ती बेरीजक्रिया असते. $\num{n+m}$ हे उदाहरणही यांत येते:
  तो $n+m$ या संख्येला अनुरूप असलेला मानक संख्यांक आहे.
\end{explain}

\end{document}
